% =====================================================================
% 附录（可选）
%
% 以下文字与附件 3 附录页完全一致，逐字照录（含附件 3 自带的格式说明
% 段落）。章标题照录为“附    录”（中间空 4 格）。
%   附录的序号编排按附录 A、附录 B……编排，
%   附录内顺序可按 B2、B2.1、B2.1.1、B2.1.2 的规律编排，
%   图表按图 B1、图 B2、表 B1、表 B2 的规律编排，与正文连续编页码。
% =====================================================================

% 正文标题“附\qquad 录”（中间空 4 格）；目录条目用不带空隙的“附录”。
\ujsunnumberedchapter{附\qquad 录}[附录]

附录（可选）主要包括一些不宜放在正文中的支撑材料，如某些原始数据、数学推导、计算程序、框图、结构图、注释、统计表、计算机打印输出件等。

“附录”二字用黑体小二号，居中，中间空4格，段前间距为2行，段后间距为1行。内容中文用宋体小四号，1.5倍行距，首行缩进2字符。

附录的序号编排按附录A、附录B．．．编排，附录（以附录B为例）内的顺序可按B2，B2.1,B2.1.1,B2.1.2的规律编排。图表按：图B1，图B2，表B1，表B2的规律编排。附录与正文装订在一起，连续编页码。

% !TEX root = ../main.tex
% 附录图表公式排版测试，保留原附录说明。
\ujsappendix
\ujschapter{附录 A\quad 图表与公式测试}

本附录检查图~\ref{fig:appendix-a1}、表~\ref{tab:appendix-a1}和式\eqref{eq:appendix-a1}。编号应从 A1 开始，与正文的图~\ref{fig:test-blue}、表~\ref{tab:test-numeric}及式\eqref{eq:test-ohm}区分。

\begin{figure}[htbp]
  \centering
  \includegraphics[width=0.38\textwidth]{test-blue.pdf}
  \caption{附录蓝色图片示例}
  \label{fig:appendix-a1}
\end{figure}

\begin{table}[htbp]
  \centering
  \caption{附录 A 模拟参数}
  \label{tab:appendix-a1}
  {\songti\zihao{-4}
  \begin{tabular}{lrr}
    \toprule
    参数组 & 电阻 / $\Omega$ & 电压 / V \\
    \midrule
    A & 10.00 & 12.50 \\
    B & 15.00 &  9.75 \\
    \bottomrule
  \end{tabular}\par}
\end{table}

附录单行公式如下：
\begin{equation}
  P = U I
  \label{eq:appendix-a1}
\end{equation}
同一附录的下一条多行公式应连续编号：
\begin{align}
  P &= I(IR) \notag\\
    &= I^2R.
  \label{eq:appendix-a2}
\end{align}
再次引用图~\ref{fig:appendix-a1}、表~\ref{tab:appendix-a1}、式\eqref{eq:appendix-a1}和式\eqref{eq:appendix-a2}，编号应保持一致。

\ujschapter{附录 B\quad 编号重置测试}

进入附录 B 后，图~\ref{fig:appendix-b1}、表~\ref{tab:appendix-b1}及式\eqref{eq:appendix-b1}均应从 B1 开始。上一附录的式\eqref{eq:appendix-a2}和正文的式\eqref{eq:test-mean}仍应跳转到原位置。

\begin{figure}[htbp]
  \centering
  \includegraphics[width=0.38\textwidth]{test-orange.pdf}
  \caption{附录橙色图片示例}
  \label{fig:appendix-b1}
\end{figure}

\begin{table}[htbp]
  \centering
  \caption{附录 B 模拟参数}
  \label{tab:appendix-b1}
  {\songti\zihao{-4}
  \begin{tabular}{lrr}
    \toprule
    参数组 & 时间 / s & 电流 / A \\
    \midrule
    C & 2.10 & 0.80 \\
    D & 3.40 & 1.25 \\
    \bottomrule
  \end{tabular}\par}
\end{table}

\begin{equation}
  E = P t
  \label{eq:appendix-b1}
\end{equation}

引用汇总：图~\ref{fig:appendix-b1}、表~\ref{tab:appendix-b1}和式\eqref{eq:appendix-b1}。附录页码与正文连续，不重新从 1 开始。

